% !TEX program = xelatex

\documentclass{article}
\usepackage{geometry}
\geometry{margin=2cm}
\usepackage{array}
\usepackage{booktabs}
\usepackage{xcolor}
\usepackage{tipa}
\usepackage{wasysym}
\usepackage{fontspec}
\usepackage{arabxetex}

\newfontfamily\arabicfont[Script=Arabic]{Amiri}
\newcommand{\ar}[1]{{\arabicfont #1}}

\newfontfamily\cursivefont{Didot}
\newcommand{\cursive}[1]{{\cursivefont #1}}

\begin{document}

\begin{center}

\cursive{\Large Additional combinations and sign:}

\vspace{1em}

\begin{tabular}{r c l}
\\
*lam-\textit{'alif}    & \ar{لا} & l\={a}       \\[6pt]
t\={a} marb\={u}\d{t}a & \ar{ة}  & -at-         \\[6pt]
hamza                  & \ar{ء}  & '            \\[6pt]
\v{s}adda              & \ar{ـّ} & {[doubling]} \\[6pt]
\textit{'alif}-madda   & \ar{آ}  & '\={a}       \\[6pt]
\end{tabular}

\vspace{2em}

The only two letter combination to have separate form in the alphabet is\\
Combination l\={a}m + \textit{'alif}

\vspace{3em}

\cursive{\Large Numerals:}

\vspace{1em}

Compound numerals are written like English, from left to right (365 = \ar{٣٦٥})

\vspace{1em}

\begin{tabular}{c c c c c c c c c c}
1 \ar{١} & 2 \ar{٢} & 3 \ar{٣} & 4 \ar{٤} & 5 \ar{٥} &
6 \ar{٦} & 7 \ar{٧} & 8 \ar{٨} & 9 \ar{٩} & 10 \ar{١٠} \\
\end{tabular}

\end{center}

\newpage

\begin{center}
{\Large The Vowel Signs.}

\vspace{0.5em}

6.1 The short vowels and the sign of quiescence:
\end{center}

\vspace{1em}

(1) fat\d{h}a, the sign for \textit{a}, is a short diagonal stroke placed over the consonant it follows
in pronunciation, as in \ar{كَتَبَ} \underline{kataba} and \ar{خَرَجَ} \underline{xaraja}.

\vspace{0.8em}

(2) kasra, the sign for \textit{i}, is the same diagonal stroke placed under the consonant it
follows in pronunciation, as in \ar{مِن} mina and \ar{بِهِ} \underline{bihi}.

\vspace{0.8em}

(3) \d{d}amma, the sign for \textit{u}, is a small w\={a}w placed over the consonant it follows in
pronunciation, as in \ar{كُتُبُ} \underline{kutubu} and \ar{رَجُلُ} \underline{rajulu}.

\vspace{0.8em}

(4) In fully vocalized texts such as the Koran, every consonant must be marked, hence
the existence of \underline{suk\={u}n}, the sign for no vowel at all (quiescence), usually written as a
small circle above the consonant, as in \ar{كَتَبْتُ} \underline{katabtu} and \ar{مِنْ} min.

\vspace{1.5em}

\begin{center}
The \textbf{long vowel} signs are as follows:
\end{center}

\vspace{0.5em}

(1) \={a} is indicated by fat\d{h}a plus alif, as in \ar{كَاتَبَا} k\={a}tab\={a} and \ar{قَامَ} q\={a}ma.
Note that \={a} is often, especially in the Koran, written defectively as ``dagger alif'' above
the consonant, as in \ar{اللّٰه} all\={a}hu and \ar{إِبْرَٰهِيمُ} 'ibr\={a}h\={i}mu.

\vspace{0.8em}

(2) \={i} is indicated by kasra plus y\={a}', as in \ar{كَبِير} kab\={i}r- and \ar{دِين} d\={i}n-.

\vspace{0.8em}

(3) \={u} is indicated by \d{d}amma plus w\={a}w, as in \ar{رَسُول} ras\={u}l- and \ar{ثُوم} \texttheta\={u}m-.

\vspace{1.5em}

6.3 The diphthong signs are a combination of the short vowel \d{a} and consonant.

\vspace{1em}

\begin{center}
(1) ay is indicated by fat\d{h}a plus y\={a}', as in \ar{أَيْنَ} 'ayna

\vspace{0.5em}

(2) aw is indicated by fat\d{h}a plus w\={a}w, as in \ar{دَوْر} dawr-
\end{center}


\vspace{1.5em}

\textbf{Otiose alif}, in certain conjugation forms an alif is appended to a lengthening w\={a}w, as in
\ar{كَتَبُوا} katab\={u}. This alif is not pronounced and serves merely to indicate the verbal form.
It owes its existence to early orthographic conventions.

\vspace{1.5em}

\begin{center}
Alif maq\d{s}\={u}ra, the alif maq\d{s}\={u}ra, also called alif bi-\d{s}\={u}rati l-y\={a}'
(Alif masquerading as y\={a}), occurs word-finally only. Written like a y\={a}', it is pronounced
exactly like a lengthening alif, as in \ar{المعنى} al-ma'n\={a} and \ar{رمى} ram\={a}. When any enclitic
suffix is added to alif bi-\d{s}\={u}rati l-y\={a}' it becomes ``tall'' alif, as in \ar{معناه} ma'n\={a}-hu and \ar{رماه}
ram\={a}-hu.
\end{center}

\vspace{1.5em}

\begin{center}
Additional Orthographic signs:
\end{center}

\vspace{0.8em}

Hamza, the sign of the glottal stop('). Word-initially it is invariably written on alif.
When the vowel of the hamza is a or u, the hamza is commonly written above the alif,
as in \ar{أرض} '\underline{ar\d{d}}- and \ar{أن} 'an.

\vspace{0.8em}

But when the vowel is I, the hamza is commonly written beneath the alif, as in \ar{إنسان}
'ins\={a}n- and \ar{إن} 'in.

\vspace{1em}

\begin{center}
Non-initially the ``bearer'' of the hamza may be:

\vspace{0.5em}

(1) alif, as in \ar{سال} sa'ala [to flow or to drip/run]

\vspace{0.5em}

(2) w\={a}w, as in \ar{سؤال} su'\={a}l- [question, inquiry/query, request]

\vspace{0.5em}

(3) y\={a} without dots, as in \ar{رئس} ra'\={i}s- [head]

\vspace{0.5em}

(4) nothing, as in \ar{نساء} nis\={a}'- [women]
\end{center}


\vspace{1.5em}

Wa\d{s}la, a small initial \d{s}\={a}d, is the sign of elision. Many initial vowels, notably the vowel
of the definite article, are elided when not in sentence-initial position. When such
elision occurs, the Wa\d{s}la is placed over the alif. E.g., when sentence initial, \ar{الأرض} 'al-'ar\d{d}u but \ar{فى الأرض} fi l-'ar\d{d}i.

\vspace{1.5em}

\v{S}adda, the sign of gemination. Doubled consonants are never written twice in Arabic
but are indicated by placing the sign \underline{\v{s}adda} over the doubled consonant. In unvocalized
texts the \underline{\v{s}adda} may be indicated sporadically, but it is not normally given.

\vspace{1em}

\begin{center}
\begin{tabular}{r l r l}
Jannat- \ar{جنّة}  & & makkat- \ar{مكّة}          \\[6pt]
-sayyid \ar{سيّد}  & nab\={i}y- (nabiyya-) \ar{نبيّ} \\[6pt]
radda \ar{ر دّ}    & nub\={u}wat-(nubuwwat) \ar{نبوّة} \\[6pt]
\end{tabular}
\end{center}

\end{document}

